% chapters/frontmatter.tex
% 前置部分: 献词 / 前言 / 使用说明 / 致谢
% 源: 原书 PDF 第 5--8 页 (印刷页 v--ix)

\thispagestyle{empty}
\vspace*{0.22\textheight}
\begin{center}
	\large\itshape
	献给 Felix Browder\\
	——我的良师与挚友,\\
	是他教会我以泛函分析学者的眼光\\
	欣赏偏微分方程
\end{center}
\cleardoublepage

\chapter*{译者序}
\markboth{译者序}{译者序}
\addcontentsline{toc}{chapter}{译者序}
\enlargethispage{2\baselineskip}

海姆·布雷齐斯 (Haïm Brezis, 1944--2024) 教授是当代偏微分方程与非线性泛函分析大师.
本书由其在巴黎大学与罗格斯大学讲授的讲义凝练而成,
将抽象泛函分析与现代 PDE (尤其是 Sobolev 空间与变分法) 浑然融合, 是享誉全球的经典教材.
为方便国内学子研读, 我们依据 Springer 2011 年英文版进行了全本中文重排与翻译.
针对电子化研读体验, 本次重排特别做了如下现代化改进与体例说明:

\begin{enumerate}[label=\arabic*., leftmargin=2.2em, itemsep=0.15em, topsep=0.15em]
	\item \textbf{习题与解答的双向交互跳转}:
	      本书收录了极其丰富的习题与综合性问题.为了让读者在自学和演算时获得最佳阅读体验,
	      我们在 PDF 电子版中全面实现了\emph{双向可逆超链接} (Bidirectional Hyperlinks):
	      \begin{itemize}[label=$\circ$, leftmargin=1.5em, itemsep=0.05em, topsep=0.05em]
		      \item 正文各章末尾的习题中, 凡在书末「部分习题解答」中附有解答的题目,
		            其方框题号 (如 \fbox{\small\bfseries 1.1}) 均内置交互跳转链接, 点击即可直接跳至书末对应解答;
		            在书末解答处, 点击加粗题号同样可瞬间跳回正文原题所在页面.
		      \item 在后置的 51 道大型「问题」与「部分问题解答」之间,
		            亦同样支持点击问题标题实现原题与参考解答之间的双向无缝跳转.
	      \end{itemize}
	\item \textbf{参考文献引用体例说明}:
	      原书参考文献列表严格遵循欧洲古典分析专著的编排传统:
	      文献按作者姓氏首字母升序排列, 但\emph{每位作者名下的著作各自独立按 $[1],[2],\dots$ 重新编号} (而非全书统一的单一大流水号).
	      因此, 正文中出现的引用形式为 ``作者名 [编号]'' (例如 ``H.~L.~Royden [1]'', ``E.~Lieb--M.~Loss [1]'', ``J.~Dieudonné [1]'').
	      文中的每一个 $[1]$ 或 $[2]$ 均已建立精确的超链接锚点, 点击即可直达书末该作者名下的具体出版条目.
	\item \textbf{图表矢量重绘与高精度交叉引用}:
	      原书中全部 9 幅几何原理示意图均采用 TikZ 矢量代码重新精细绘制, 避免了扫描位图的失真与模糊;
	      定理、引理、推论、小节编号 ($\S$) 以及公式引用均全面绑定超链接 (\verb|\ref| 与 \verb|\eqref|),
	      确保读者可快速穿透引用链条.
	\item \textbf{翻译原则与严谨性}:
	      本译本坚持严谨忠实原则, 绝不擅自删节、合并或跳跃原文论证细节.
	      对于原版英文印刷中已发现的个别符号笔误, 译本均在 LaTeX 源文件中以注释形式详实记录并予以严格核对.
\end{enumerate}

\vspace{0.6\baselineskip}
\begin{flushright}
	译者\\
	2026 年 10 月
\end{flushright}

\cleardoublepage

\chapter*{前言}
\markboth{前言}{前言}
\addcontentsline{toc}{chapter}{前言}

本书源于我在巴黎大学多年讲授的一门课程.它面向在实分析方面已有良好基础的学生
(例如 G.\,B.~Folland [2], A.\,W.~Knapp [1] 和 H.\,L.~Royden [1] 的教科书中所讲述的内容).
我构想了一套混合了两个不同``世界''之内容的方案: 泛函分析 (FA) 与偏微分方程 (PDE).
第一部分处理 FA 与算子理论中的抽象结果.第二部分研究具有特定可微性质的函数空间
(一个或多个实变量的函数), 即著名的 Sobolev 空间, 它位于现代 PDE 理论的核心.
我说明如何把 FA 中的抽象结果用于求解 PDE.Sobolev 空间出现在纯数学与应用数学的广泛问题中.
它们出现在例如微分几何, 调和分析, 工程, 力学和物理中产生的线性与非线性 PDE 里.
它们属于任何分析方向研究生的工具箱.

遗憾的是, FA 与 PDE 虽然联系紧密, 却常常分属两门课程.FA 中处理的许多问题都源自 PDE
(从历史的角度看, 可参见例如 J.~Dieudonné [1] 与 H.~Brezis--F.~Browder [1]).
专论 FA 的书 (甚至多卷本巨著) 数量众多.处理 PDE 的教科书也不少.
然而, 面向研究生的综合性讲述却很少见, 我试图填补这一空缺.
常常被数学中最抽象的构造所吸引的学生, 通常也为 FA 的优雅所吸引.
另一方面, 他们又对 PDE 中没完没了的公式以及数不清的下标感到畏惧.
我尝试先分析一维 PDE (即 ODE——常微分方程) 这一简单情形, 从而呈现一条从 FA 到 PDE 的
``平滑''过渡; 对初学者而言, 一维情形看起来要容易把握得多.按这种做法, 我所阐述的技巧
对 ODE 而言可能过于精细, 但它们后来成为 PDE 理论的基石.
这种安排使学生此后处理精细的高维 PDE 时容易得多.

本书的一个较早版本最初于 1983 年以法文出版, 此后有多种译本, 在世界范围内颇受欢迎,
并被许多欧洲大学采用为教科书.法文本的一个缺陷是缺少习题.
本书含有丰富的习题.我打算在未来的版本中再增加一些.
我也概述了一些较新的进展, 特别是在非线性 PDE 方向上的进展.

\section*{使用说明}
\addcontentsline{toc}{section}{使用说明}

\begin{enumerate}[label=\arabic*., leftmargin=2.2em, itemsep=0.35em]
	\item 以圆点符号 \textbullet\ 起头的陈述或段落极为重要, 务必透彻掌握, 才能理解后文.
	\item 标有星号 $\star$ 的结果初学者可以跳过; 它们只对高年级读者有意义.
	\item 在每一章中, 我把命题, 定理与推论连续编号 (例如命题~
\ref{prop:3.6} 之后是定理~
\ref{thm:3.7}, 推论~
\ref{cor:3.8},
	      等等).只有注记与引理另行单独编号.
	\item 为简化表述, 我假定所有向量空间都是 $\mathbb{R}$ 上的.多数结果对 $\mathbb{C}$ 上的向量空间
	      仍然成立.我在第~
\ref{chap:11}~章加了一节简短的讨论, 说明其异同.
	\item 许多章之后附有大量习题.部分解答放在书末.更精细的问题另设一节``问题'',
	      其后是``问题的部分解答''.这些问题通常需要用到若干章的知识.
	      我在每个问题的开头标明了所涉及的章.有些习题与问题阐述的是本章正文中未加细节
	      或未给证明的结果.
\end{enumerate}

\section*{致谢}
\addcontentsline{toc}{section}{致谢}

在本书的写作过程中, 我得到了两位挚友兼前同事的诸多鼓励: Ph.~Ciarlet 与 H.~Berestycki.
我非常感谢 G.~Tronel, M.~Comte, Th.~Gallouet, S.~Guerre-Delabrière, O.~Kavian, S.~Kichenassamy 以及已故的 Th.~Lachand-Robert, 他们与我分享了与学生打交道的``实战经验''.
S.~Antman, D.~Kinderlehrer 和 Y.~Li 向我解释了美国学生的背景与``品味''.
C.~Jones 惠允我使用他为个人之需所准备的, 法文原书若干章的英译稿.
我要感谢 A.~Ponce, H.-M.~Nguyen, H.~Castro 和 H.~Wang, 他们仔细校阅了本书的部分内容.
我有幸在罗格斯大学得到两位出色的助手, 本书的大部分内容由她们录入:
现已退休的 Barbara Miller, 以及现任的 Barbara Mastrian.对她们始终如一的奉献与专业帮助,
我的赞美与感激之词不足以表达.她们总能为 PDE 公式中棘手的繁难之处找到令人满意的处理办法.
没有她们的热情与耐心, 本书绝无可能完成.一如既往, 与 Springer 的 Ann Kostant 在这一项目上的合作
令我十分愉快.过去我多次有机会体会到她对数学界长期不渝的投入.

作者部分得到 NSF 基金 DMS-0802958 的资助.

\vspace{1.5\baselineskip}
\begin{flushright}
	Haim Brezis\\
	罗格斯大学\\
	2010 年 3 月
\end{flushright}

\cleardoublepage
